\section{Integrated Discussion}

\subsection{What is strongest in the repository}

The full repository does not support every attractive story equally well. Its strongest supported thesis has three parts. First, there is a real shared difficulty axis across humans and machine systems on ARC. Second, validated final solver structure is one of the strongest explanatory variables for LLM difficulty in the entire repo. Third, human burden is not captured by that same structure to the same degree and instead looks more tied to time cost, search burden, and pair-level heterogeneity.

\subsection{What should be treated as central versus secondary}

If the paper is trying to be intellectually honest, it should rank its findings rather than narrate them all at the same volume. The central retained claims are: ARC looks psychometrically clean on the LLM side; sparse human ARC data are usable once reliability is estimated explicitly; human-versus-average-model alignment is real but bounded by human split-half reliability; validated solver structure predicts LLM difficulty strongly; the solver-structure link is significantly stronger for LLM difficulty than for human difficulty; and current non-LLM systems do not reproduce the human difficulty profile as well as the LLM consensus does. Those are the pieces the repo keeps recovering from multiple directions.

Secondary but still worthwhile claims include the closeness augmentation result, which suggests softer output similarity adds human-relevant signal for duration and latent ease; the pooled cross-ARC extension, which supports the main asymmetry on a wider benchmark footprint; and the non-LLM ``middle'' pattern in the complexity panel, which is suggestive but underpowered. The paper should treat these as useful amplifiers and qualifiers, not as the main pillars of the argument.

\subsection{How the non-LLM result should be narrated}

The non-LLM section is likely to be the hardest part of the final paper to get right rhetorically. It will be tempting either to bury it because it is weaker than hoped or to oversell its complementarity because the raw psychometric result is disappointing. Neither move would help. The most credible framing is this: current non-LLM systems are not yet the clean human-like alternative to the LLM consensus on ARC, but they are also not vacuous. They contribute complementary successes, expose a split between benchmark score and human-like item placement, and converge architecturally on a loop that may be important even where the current psychometric results remain modest.

\subsection{What the paper should explicitly flag as pressure points}

Several pressure points deserve explicit flagging in the manuscript rather than quiet footnotes. The first is the tiny fully matched human-plus-LLM-plus-non-LLM-plus-complexity overlap, which is effectively an ARC-2 comparison on 17 independent tasks. The second is the task-level versus pair-level mismatch: final solver files are task objects, but human burden often varies by pair. The third is the sampling character of the human dataset, which is sparse, opportunistic, and not a balanced exam design. The fourth is benchmark drift between ARC-1 and ARC-2, including the fact that the six shared eval IDs preserve test examples but change training examples. The fifth is that final approved solver size is not minimal solver size, still less human conceptual complexity.

There is also a more subtle pressure point: exact-match success is both the right leaderboard metric and a lossy behavioral summary. The solution-closeness addenda show that softer metrics do carry some residual signal, especially for human time cost, but they do not overturn the exact-match complexity story. The paper should say this clearly, because otherwise readers may assume either that binary scoring is the whole truth or that any soft metric automatically invalidates it. Neither reading is right.

\subsection{What the paper should not claim}

The manuscript should avoid at least three overclaims. It should not say that current LLMs and humans solve ARC for the same reasons. The human-versus-LLM complexity asymmetry is direct evidence against that. It should not say that current non-LLM systems are broadly more human-like than LLMs on the available psychometric evidence. They are not. And it should not keep the original thinking-advantage result as a clean flagship claim. After label verification and floor-sensitive checks, that result is better described as an audited lead that did not hold up cleanly.

\subsection{What future work would actually matter}

The future-work section should not devolve into generic calls for more data. The repository already points to the highest-value next steps: increase the fully matched overlap, make the task-level versus pair-level distinction explicit across the manuscript, carry the strongest closeness sensitivity analyses forward without displacing the main binary complexity result, expand cross-ARC linkage carefully, and add a small number of high-yield figures that visualize the decomposition story directly.

More specifically, the next empirical milestone would be to obtain a larger human benchmark with richer pair-level and response-grid logging so that solver structure, exact correctness, wrong-answer closeness, and duration can all be modeled on the same tasks. On the machine side, the most valuable next step is probably not more leaderboard chasing. It is to test whether the architectural loop suggested by URM, TRM, VARC, and McGovern can be made to yield a more human-like item profile, or whether those systems will continue to improve score while drifting away from the human ordering the way TRM seems to in the current repo.

\subsection{What this first full-paper draft is actually for}

As a first full-paper draft, the main job of this manuscript is not to finish every figure or settle every caveat. It is to make the repository legible as one coherent scientific argument. The argument it currently supports is strong enough to be worth taking seriously and narrow enough to be scrutinized. ARC contains a real shared difficulty axis across humans and machine systems, but it does not impose identical burden on all of them. LLM difficulty is strongly linked to the structure of validated successful solvers. Human difficulty is linked more weakly to that structure and more strongly to time, search, and pair heterogeneity. Current non-LLM systems remain the most promising architectural alternative and the weakest psychometric match. That is a sufficiently sharp thesis to organize a much larger final manuscript around.
